\section{Conclusion and Future Work}
\label{sec:conclusion}

We presented \texttt{Rust2Isabelle}, an Isabelle/HOL backend for the Charon--Aeneas toolchain.
It combines Aeneas's functional translation with Isabelle-specific representations and primitives' support.
Our evaluation demonstrates the feasibility of the approach and delineates its current coverage.

\paragraph{Limitations.}
The prototype inherits Aeneas's restrictions on unsafe code and interior mutability.
Floating-point operations and some Rust library built-ins remain unsupported, and the backend does not explicitly represent computations that diverge.
Generated theories may still use \texttt{axiomatization} for external declarations without bodies and unsupported operations.
The translation does not yet have a mechanized end-to-end semantics-preservation proof, and the backend does not generate functional-correctness specifications or proofs.


\paragraph{Future work.}
Future work will broaden the accepted Rust fragment by adding floating-point types and definitions for more Rust library built-ins to the \rprelude{}.
An explicit representation of divergence is also needed for computations that may not terminate.

Longer-term work can use translation validation to establish semantic preservation by generating Isabelle/HOL proof obligations for each Pure-to-Isabelle translation, rather than proving the entire translator correct~\cite{Pnueli1998}.
Rust specifications could also be translated into Isabelle/HOL statements for functional-correctness proofs.
Supporting unsafe code and interior mutability is more challenging, because it would require an explicit memory model.
